\begin{thebibliography}{99}
\small
\bibitem{P03} Cohn, C., et al. (n.d.). A theory of adaptive scaffolding for LLM-based pedagogical agents.
\bibitem{P04} LearnLM Team, Google and Eedi. (2025). AI tutoring can safely and effectively support students: An exploratory RCT in UK classrooms.
\bibitem{P06} Jang, H., and So, H.-J. (n.d.). Adaptive scaffolding with generative AI chatbot for solving ill-structured problems.
\bibitem{P07} Chi, M., VanLehn, K., Litman, D., and Jordan, P. (2011). An evaluation of pedagogical tutorial tactics for a natural language tutoring system: A reinforcement learning approach. \emph{International Journal of Artificial Intelligence in Education}, 21, 83--113. doi:10.3233/JAI-2011-014.
\bibitem{P08} Weinberger, A., and Fischer, F. (n.d.). A framework to analyze argumentative knowledge construction in computer-supported collaborative learning.
\bibitem{P10} Chai, W. N., et al. (n.d.). Beyond answer generation: Scaffolding metacognitive thinking via LLM-driven Socratic dialogue.
\bibitem{P11} Jin, F. J.-Y., et al. (n.d.). Characterising student behaviours with AI-generated feedback in higher education.
\bibitem{P12} Barcaui, A. (n.d.). ChatGPT as a cognitive crutch: Evidence from a randomized controlled trial on knowledge retention.
\bibitem{P13} de Jong, S., Paananen, V., Tag, B., and van Berkel, N. (n.d.). Cognitive forcing for better decision-making: Reducing overreliance on AI systems through partial explanations. doi:10.1145/3710946.
\bibitem{P14} VanLehn, K., et al. (2007). When are tutorial dialogues more effective than reading? \emph{Cognitive Science}.
\bibitem{P15} Misiejuk, K., L\'opez-Pernas, S., Kaliisa, R., and Saqr, M. (n.d.). Cognitive offloading in student--AI collaboration: A longitudinal analysis of prompting strategies.
\bibitem{P18} Yin, Z., et al. (n.d.). Do large language models know what they don't know?
\bibitem{P22} Ackerman, C. (2026). Evidence for limited metacognition in LLMs. \emph{International Conference on Learning Representations}.
\bibitem{P23} Cohn, C., et al. (n.d.). Evidence--Decision--Feedback: Theory-driven adaptive scaffolding for LLM agents.
\bibitem{P25} Lee, H.-Y., Chen, P.-H., Huang, Y.-M., and Wu, T.-T. (2026). From immediate answers to structured inquiry: A Socratic method-enhanced ChatGPT system for blended learning discussions. \emph{Journal of Educational Computing Research}, 64(5), 1254--1293. doi:10.1177/07356331261429306.
\bibitem{P29} Ilkou, E., Linzbach, S., and Wallat, J. (n.d.). Hybrid evaluation of Socratic dialogue for teaching.
\bibitem{P32} Hashemi, H., Eisner, J., Rosset, C., Van Durme, B., and Kedzie, C. (2024). LLM-RUBRIC: A multidimensional, calibrated approach to automated evaluation of natural language texts. \emph{Proceedings of the 62nd Annual Meeting of the Association for Computational Linguistics}, 13806--13834.
\bibitem{P33} Cui, J., Pu, K., and Grossman, T. (n.d.). LOOM: Personalized learning informed by daily LLM conversations toward long-term mastery via a dynamic learner memory graph.
\bibitem{P34} Griot, M., Hemptinne, C., Vanderdonckt, J., and Yuksel, D. (n.d.). Large language models lack essential metacognition for reliable medical reasoning. doi:10.1038/s41467-024-55628-6.
\bibitem{P35} LearnLM Team, Google. (2024). LearnLM: Improving Gemini for learning.
\bibitem{P38} Zhu, Q., Li, X., Dong, Y., Chang, P., and Fan, M. (2026). Not all cognitive offloading is equal: Distinguishing dependent and autonomous offloading to generative AI. \emph{Frontiers in Psychology}, 17, 1878629. doi:10.3389/fpsyg.2026.1878629.
\bibitem{P39} Neagu, A., Wong, J. T. H., Messer, M., Nelson, R., and Johnson, P. B. (n.d.). Rethinking scaffolding in LLM tutors: The interactional mismatch between benchmarks and real-world deployments.
\bibitem{P41} Ma, T., Chen, R., Li, A. T., and Liu, H. (2024). Socratic ChatGPT: Theory, design, and empirical evaluations. \emph{Pacific Asia Conference on Information Systems Proceedings}.
\bibitem{P42} Chi, M. T. H., and Wylie, R. (n.d.). The ICAP framework: Linking cognitive engagement to active learning outcomes.
\bibitem{P43} Lee, H.-P., Sarkar, A., Tankelevitch, L., Drosos, I., Rintel, S., Banks, R., and Wilson, N. (n.d.). The impact of generative AI on critical thinking: Self-reported reductions in cognitive effort and confidence effects from a survey of knowledge workers. doi:10.1145/3706598.3713778.
\bibitem{P45} Zhai, C., Wibowo, S., and Li, L. D. (2024). The effects of over-reliance on AI dialogue systems on students' cognitive abilities: A systematic review. \emph{Smart Learning Environments}, 11, 28. doi:10.1186/s40561-024-00316-7.
\bibitem{P46} Park, S., Subramonyam, H., and Kulkarni, C. (n.d.). Thinking assistants: LLM-based conversational assistants that help users think by asking rather than answering.
\bibitem{P49} Jurenka, I., et al. (2025). Towards responsible development of generative AI for education: An evaluation-driven approach.
\bibitem{P50} Maurya, K. K., Srivatsa, K. V. A., Petukhova, K., and Kochmar, E. (2025). Unifying AI tutor evaluation: An evaluation taxonomy for pedagogical ability assessment of LLM-powered AI tutors. \emph{Proceedings of NAACL-HLT}, 1234--1251.
\bibitem{P52} Zhang, A., et al. (n.d.). When help is unhelpful: Evaluating AI tutors for productive struggle.
\bibitem{P53} Kosmyna, N., et al. (n.d.). Your brain on ChatGPT: Accumulation of cognitive debt when using an AI assistant for essay writing task. Preprint.
\bibitem{P54} Gonsalves, C. (2026). Generative AI's impact on critical thinking: Revisiting Bloom's taxonomy. \emph{Journal of Marketing Education}, 48(1), 4--19. doi:10.1177/02734753241305980.
\end{thebibliography}
